\subsection{Sheaves of sets}

Let B be a topological space.

DEFINITION 1. --- A presheaf on B, relative to a base \(\mathcal{B}\) of the topology of B, is a projective system of sets, relative to the index set \(\mathcal{B}\) ordered by the inclusion relation.

In other words (E, III, p.~52), a presheaf \(\mathcal{F}\) on B relative to \(\mathcal{B}\) is a pair \(((\mathcal{F}(\mathrm{U}))_{\mathrm{U}\in\mathcal{B}},(f_{\mathrm{UV}}))\) , also denoted \((\mathcal{F}(\mathrm{U}),f_{\mathrm{UV}})\) , where \((\mathcal{F}(\mathrm{U}))_{\mathrm{U}\in\mathcal{B}}\) is a family of sets with \(\mathcal{B}\) as index set, and where for each pair (U,V) of elements of \(\mathcal{B}\) such that U \(\subset\) V, \(f_{\mathrm{UV}}\) is a map from \(\mathcal{F}(\mathrm{V})\) into \(\mathcal{F}(\mathrm{U})\) ; these maps satisfy the following conditions:

\((\mathrm{PF}_{1})\) The relations \(\mathrm{U}\subset \mathrm{V}\subset \mathrm{W}\) imply \(f_{\mathrm{UW}} = f_{\mathrm{UV}}\circ f_{\mathrm{VW}}\).

\((\mathrm{PF}_2)\) For every open set \(\mathbf{U} \in \mathcal{B}\) , \(f_{\mathrm{UU}}\) is the identity map of \(\mathcal{F}(\mathbf{U})\) .

A presheaf on B relative to the set of open subsets of B is simply called a presheaf on B.

Let \(\mathcal{F} = (\mathcal{F}(\mathrm{U}), f_{\mathrm{UV}})\) be a presheaf on B, relative to a base \(\mathcal{B}\) of the topology of B. Let U be an element of the base \(\mathcal{B}\) . The elements of \(\mathcal{F}(\mathrm{U})\) are called the sections of \(\mathcal{F}\) on U. If V is an element of the basis \(\mathcal{B}\) containing U and \(s\) is an element of \(\mathcal{F}(\mathrm{V})\) , the element \(f_{\mathrm{UV}}(s)\) of \(\mathcal{F}(\mathrm{U})\) is called the restriction of \(s\) to U. If there is no risk of ambiguity regarding the mappings \(f_{\mathrm{UV}}\) ; we shall denote by \(s \mid \mathrm{U}\) the restriction of \(s\) to U.

Let B′ be an open subset of B, and \(\mathcal{B}'\) be a base of the topology of B' such that \(\mathcal{B}' \subset \mathcal{B}\). We call the presheaf on B', relative to \(\mathcal{B}'\) deduced from \(\mathcal{F}\) by restriction, and we denote by \(\mathcal{F} \mid \mathcal{B}'\) the projective system \(((\mathcal{F}(U))_{U \in \mathcal{B}'}\) , \((f_{UV})\) deduced from \(\mathcal{F}\) by restriction to \(\mathcal{B}'\) of the index set (loc. cit.). When \(\mathcal{F}\) is a presheaf on B and \(\mathcal{B}'\) is the set of open subsets of B', the presheaf \(\mathcal{F} \mid \mathcal{B}'\) is also denoted \(\mathcal{F} \mid B'\) and called the presheaf deduced from \(\mathcal{F}\) by restriction to B'.

DEFINITION 2. --- Let \(\mathcal{F} = (\mathcal{F}(U), f_{UV})\) be a presheaf on B. We say that \(\mathcal{F}\) is a sheaf on B if, for every open subset U of B and every family \((U_i)_{i \in I}\) of open subsets of B, with union U, the following properties are satisfied:

\((\mathrm{F}_{1})\) The map \((f_{\mathrm{U}_{i}\mathrm{U}})_{i\in\mathrm{I}}\colon\mathcal{F}(\mathrm{U})\to\prod_{i\in\mathrm{I}}\mathcal{F}(\mathrm{U}_{i})\) is injective;

\((\mathrm{F}_{2})\) For every family \((s_{i}) \in \prod_{i \in \mathrm{I}} \mathcal{F}(\mathrm{U}_{i})\) such that \(f_{(\mathrm{U}_{i} \cap \mathrm{U}_{j})\mathrm{U}_{i}}(s_{i}) = f_{(\mathrm{U}_{i} \cap \mathrm{U}_{j})\mathrm{U}_{j}}(s_{j})\) for every pair \((i, j) \in \mathrm{I} \times \mathrm{I}\) , there exists an element s of \(\mathcal{F}(\mathrm{U})\) such that for every \(i \in I\), we have \(f_{\mathrm{U}_{i}\mathrm{U}}(s) = s_{i}\) .

Note. --- Let \(\mathcal{F}\) be a presheaf on B. For every open set U of B, \(f_{\emptyset \mathrm{U}}\) is a map from \(\mathcal{F}(\mathrm{U})\) into \(\mathcal{F}(\emptyset)\) , therefore \(\mathcal{F}(\emptyset)\) is nonempty as soon as there exists an open set U for which \(\mathcal{F}(\mathrm{U})\) is nonempty. If \(\mathcal{F}\) is a sheaf, \(\mathcal{F}(\emptyset)\) is a singleton set; this is seen by applying \((\mathrm{F}_{1})\) and \((\mathrm{F}_{2})\) to the covering of the empty set by the empty family \((\mathrm{I}=\varnothing)\) .

Let F be a sheaf on B and let B' be an open subset of B; the presheaf \(F|B'\) deduced from F by restriction to B' is a sheaf, called the sheaf deduced from F by restriction to B'.

